\documentclass[10pt,a4paper]{amsart}
\usepackage{amsmath,amssymb,amsthm,graphicx}
\usepackage{verbatim}
\newenvironment{reference}{\comment}{\endcomment}
%\newenvironment{reference}{}{}
\newenvironment{slogan}{\comment}{\endcomment}
\newenvironment{history}{\comment}{\endcomment}

% For commutative diagrams we use Xy-pic
\usepackage[all]{xy}

% We use 2cell for 2-commutative diagrams.
\xyoption{2cell}
% Theorem environments.
%
\theoremstyle{plain}
\newtheorem{theorem}[subsection]{Theorem}
\newtheorem{proposition}[subsection]{Proposition}
\newtheorem{lemma}[subsection]{Lemma}

\theoremstyle{definition}
\newtheorem{definition}[subsection]{Definition}
\newtheorem{example}[subsection]{Example}
\newtheorem{exercise}[subsection]{Exercise}
\newtheorem{situation}[subsection]{Situation}

\theoremstyle{remark}
\newtheorem{remark}[subsection]{Remark}
\newtheorem{remarks}[subsection]{Remarks}

\numberwithin{equation}{subsection}

% Macros
%
\def\lim{\mathop{\mathrm{lim}}\nolimits}
\def\colim{\mathop{\mathrm{colim}}\nolimits}
\def\Spec{\mathop{\mathrm{Spec}}}
\def\Hom{\mathop{\mathrm{Hom}}\nolimits}
\def\Ext{\mathop{\mathrm{Ext}}\nolimits}
\def\SheafHom{\mathop{\mathcal{H}\!\mathit{om}}\nolimits}
\def\SheafExt{\mathop{\mathcal{E}\!\mathit{xt}}\nolimits}
\def\Sch{\mathit{Sch}}
\def\Mor{\mathop{\mathrm{Mor}}\nolimits}
\def\Ob{\mathop{\mathrm{Ob}}\nolimits}
\def\Sh{\mathop{\mathit{Sh}}\nolimits}
\def\NL{\mathop{N\!L}\nolimits}
\def\CH{\mathop{\mathrm{CH}}\nolimits}
\def\proetale{{pro\text{-}\acute{e}tale}}
\def\etale{{\acute{e}tale}}
\def\QCoh{\mathit{QCoh}}
\def\Ker{\mathop{\mathrm{Ker}}}
\def\Im{\mathop{\mathrm{Im}}}
\def\Coker{\mathop{\mathrm{Coker}}}
\def\Coim{\mathop{\mathrm{Coim}}}

% Boxtimes
%
\DeclareMathSymbol{\boxtimes}{\mathbin}{AMSa}{"02}

%
% Macros for moduli stacks/spaces
%
\def\QCohstack{\mathcal{QC}\!\mathit{oh}}
\def\Cohstack{\mathcal{C}\!\mathit{oh}}
\def\Spacesstack{\mathcal{S}\!\mathit{paces}}
\def\Quotfunctor{\mathrm{Quot}}
\def\Hilbfunctor{\mathrm{Hilb}}
\def\Curvesstack{\mathcal{C}\!\mathit{urves}}
\def\Polarizedstack{\mathcal{P}\!\mathit{olarized}}
\def\Complexesstack{\mathcal{C}\!\mathit{omplexes}}
% \Pic is the operator that assigns to X its picard group, usage \Pic(X)
% \Picardstack_{X/B} denotes the Picard stack of X over B
% \Picardfunctor_{X/B} denotes the Picard functor of X over B
\def\Pic{\mathop{\mathrm{Pic}}\nolimits}
\def\Picardstack{\mathcal{P}\!\mathit{ic}}
\def\Picardfunctor{\mathrm{Pic}}
\def\Deformationcategory{\mathcal{D}\!\mathit{ef}}
\usepackage{cleveref}
\crefname{originallemma}{original lemma}{original lemmas}
\Crefname{originallemma}{Original lemma}{Original lemmas}
\makeatletter
\newlabel{lemma-finite-type-points-surjective-morphism}{{06EJ}{source}{}{sourcepointer-1}{}}
\newlabel{lemma-finite-type-points-surjective-morphism@cref}{{[originallemma][0][]06EJ}{[1][0][]source}{}{}{}}
\makeatother
\crefname{originallemma}{original lemma}{original lemmas}
\Crefname{originallemma}{Original lemma}{Original lemmas}
\makeatletter
\newlabel{morphisms-lemma-identify-finite-type-points}{{02J2}{source}{}{sourcepointer-2}{}}
\newlabel{morphisms-lemma-identify-finite-type-points@cref}{{[originallemma][0][]02J2}{[1][0][]source}{}{}{}}
\makeatother
\crefname{originallemma}{original lemma}{original lemmas}
\Crefname{originallemma}{Original lemma}{Original lemmas}
\makeatletter
\newlabel{morphisms-lemma-jacobson-finite-type-points}{{01TB}{source}{}{sourcepointer-3}{}}
\newlabel{morphisms-lemma-jacobson-finite-type-points@cref}{{[originallemma][0][]01TB}{[1][0][]source}{}{}{}}
\makeatother
\crefname{originalremark}{original remark}{original remarks}
\Crefname{originalremark}{Original remark}{Original remarks}
\makeatletter
\newlabel{spaces-properties-remark-list-properties-local-etale-topology}{{03E7}{source}{}{sourcepointer-4}{}}
\newlabel{spaces-properties-remark-list-properties-local-etale-topology@cref}{{[originalremark][0][]03E7}{[1][0][]source}{}{}{}}
\makeatother
\title{Over a Jacobson base the finite type points have transcendence degree zero}
\author{The Stacks Project Authors}
\date{stacks-project:spaces-morphisms.tex:local-snapshot-20261012; GNU FDL 1.2 or later}
\begin{document}
\maketitle
\setcounter{subsection}{1}
\noindent\textit{Editorial note.} All delivered mathematical text is the author's own, copied from the pinned Stacks Project chapter file spaces-morphisms.tex; no statement or proof step was rewritten, reordered or removed.\par\medskip
\noindent\textit{Editorial note.} The proof is a presentation reduction that produces a closed point in the presentation and invokes the scheme case; the shrinking step is justified by the cited identification of finite type points.\par\medskip
\noindent\textit{Editorial note.} The four references of the counted spans are recorded as non-load-bearing navigation pointers; the argument that consumes them stays inside the counted proof span.\par\medskip
\noindent\textit{Editorial note.} Preamble: the three native excerpt spans of sources/stacks-spaces-morphisms/preamble.tex (lines 8-13, 14-18 and 154-217). The chapter file itself carries no preamble text.\par\medskip
\noindent\textit{Allowed standard prerequisites (editorial).}\begin{itemize}
\item Lifting finite type points along a locally finite type surjective morphism. Lemma lemma-finite-type-points-surjective-morphism (44058), used at proof lines 6824-6825. Produces the finite type point u of U over x.
\item Finite type points of a scheme can be made closed after shrinking. Morphisms, Lemma morphisms-lemma-identify-finite-type-points, used at proof lines 6826-6827. Makes u a closed point of the shrunk U.
\item The scheme case over a Jacobson base. Morphisms, Lemma morphisms-lemma-jacobson-finite-type-points, used at proof lines 6828-6830. Gives that kappa(u)/kappa(v) is finite.
\item The definition of Jacobson for an algebraic space. Properties of Spaces, Remark spaces-properties-remark-list-properties-local-etale-topology, used at proof lines 6814-6816. Fixes the hypothesis on Y.
\end{itemize}\medskip
\ifdefined\hypertarget\hypertarget{sourcepointer-1}{}\fi
\noindent\textit{Source reference.} Original lemma 06EJ: focused-seeds/AGP/sources/stacks-snapshot-20261011/tags.txt, source lines 8218--8218. Navigation pointer into the Stacks Project; the standard fact it supplies is applied inside the counted proof span and the argument resting on it stays there.\par\medskip
\ifdefined\hypertarget\hypertarget{sourcepointer-2}{}\fi
\noindent\textit{Source reference.} Original lemma 02J2: focused-seeds/AGP/sources/stacks-snapshot-20261011/tags.txt, source lines 3321--3321. Navigation pointer into the Stacks Project; the standard fact it supplies is applied inside the counted proof span and the argument resting on it stays there.\par\medskip
\ifdefined\hypertarget\hypertarget{sourcepointer-3}{}\fi
\noindent\textit{Source reference.} Original lemma 01TB: focused-seeds/AGP/sources/stacks-snapshot-20261011/tags.txt, source lines 2368--2368. Navigation pointer into the Stacks Project; the standard fact it supplies is applied inside the counted proof span and the argument resting on it stays there.\par\medskip
\ifdefined\hypertarget\hypertarget{sourcepointer-4}{}\fi
\noindent\textit{Source reference.} Original remark 03E7: focused-seeds/AGP/sources/stacks-snapshot-20261011/tags.txt, source lines 4446--4446. Navigation pointer into the Stacks Project; the standard fact it supplies is applied inside the counted proof span and the argument resting on it stays there.\par\medskip
\begin{lemma}
\label{lemma-jacobson-finite-type-points}
Let $S$ be a scheme. Let $f : X \to Y$ be a morphism of algebraic spaces
over $S$. If $f$ is locally of finite type, $Y$ is Jacobson
(Properties of Spaces, Remark
\ref{spaces-properties-remark-list-properties-local-etale-topology}),
and $x \in |X|$ is a finite type point of $X$,
then the transcendence degree of $x/f(x)$ is $0$.
\end{lemma}

\begin{proof}
Choose a scheme $V$ and a surjective \'etale morphism $V \to Y$.
Choose a scheme $U$ and a surjective \'etale morphism $U \to X \times_Y V$.
By Lemma \ref{lemma-finite-type-points-surjective-morphism}
we can find a finite type point $u \in U$ mapping to $x$.
After shrinking $U$ we may assume $u \in U$ is closed
(Morphisms, Lemma \ref{morphisms-lemma-identify-finite-type-points}).
Let $v \in V$ be the image of $u$. By
Morphisms, Lemma \ref{morphisms-lemma-jacobson-finite-type-points}
the extension $\kappa(u)/\kappa(v)$ is finite.
This finishes the proof.
\end{proof}

\end{document}
